\hnsection{起源}

近一个世纪以来被称为“李群理论”的理论，基本上由一位数学家 Sophus Lie 发展起来。

在开始叙述这段历史之前，我们先简要概述为此铺平道路的若干早期研究。

\textbf{(a) 变换群（Klein-Lie，1869--1872）}

约在 1860 年，一个有限集合的置换群理论正在发展，并开始得到应用（Serret、Kronecker、Mathieu、Jordan）。另一方面，正蓬勃发展的不变量理论使数学家熟悉了某些在复合下稳定的几何变换的无限集合（尤其是线性变换或射影变换）。然而，在 Jordan {[}7{]} 于 1868 年研究“运动群”（三维欧氏空间位移群的闭子群）之前，似乎还没有人有意识地把这两股思想联系起来。

1869 年，年轻的 Felix Klein（1849--1925）在柏林与挪威人 Sophus Lie（1842--1899）结为朋友；Lie 比他年长几岁，二人因共同关心 Plücker 的“直线几何”，尤其是直线复形理论，而走到了一起。大约就在这时，Lie 形成了他最具原创性的思想之一：把不变量的概念引入分析学和微分几何；其思想来源之一，是他注意到微分方程的经典“求积”积分方法完全依赖于方程在一个“连续”变换族下保持不变。1869 年的第一篇著作（由 Klein 编订）中，Lie 运用了这一思想；他在那里研究了“Reye 复形”（由四条直线组成的集合：这些直线在四面体的各面上截出四个具有给定交比的点）以及以该复形的直线为切线的曲线和曲面 {[}3a{]}：他的方法依赖于 Reye 复形在保持四面体顶点不变的三参数交换群（PGL(4, C) 的极大环面）下的不变性。Klein 和 Lie 于 1870 年春在巴黎共同写作的论文中也贯穿着同一思想 {[}1, a{]}；在那里，他们基本确定了平面射影群 PGL(3, C) 的交换连通子群，并研究了其轨道的几何性质（称为曲线或曲面 V）；通过统一的步骤，这使他们获得了各种代数或超越曲线的性质，例如 \(y = cx^m\) 或对数螺线。两篇论文都突出了 Galois 和 Jordan 的理论给他们留下的深刻印象（Jordan 关于 Galois 的评论于 1869 年发表在 Math. Annalen；另一方面，Lie 早在 1863 年就听说过 Galois 理论）。Klein 于 1871 年开始对非欧几何感兴趣，并由此开始研究一种对所有已知几何进行分类的原则，最终于 1872 年形成“Erlangen 纲领”。至于 Lie，他在 1873 年致 A. Mayer 的一封信中（{[}3{]}，第 V 卷，第 584 页）把自己关于变换群的思想追溯到巴黎逗留时期；在 1871 年的一篇著作中（{[}3 b{]}，第 208 页），他已经使用“变换群”一词，并明确提出确定 GL(n, C) 的所有子群（“连续或不连续”）的问题。坦率地说，Klein 和 Lie 在进入这个新的数学世界时都遇到了一些困难，Klein 把 Jordan 新近出版的“论著”说成是一本“七印封缄的书”（{[}2{]}，第 51 页）；此外，他在谈到 {[}1 a{]} 和 {[}1 b{]} 时写道：“关于连续算子群这一启发性思想的一切功劳都应归于 Lie，特别是微分方程和偏导数的积分。后来他在连续群理论中发展出的所有概念，虽然都已以萌芽形式存在，但阐述得如此不充分，以至于只有经过长时间交谈，我才能使他相信许多细节，例如首先要相信曲线 V 的确存在”（{[}2{]}，第 415 页）。

\textbf{(b) 无穷小变换}

“无穷小”变换的概念至少可以追溯到无穷小分析学的开端；我们知道，笛卡尔承认“在无穷小范围内”每个平面运动都可视为旋转时，发现了瞬时旋转中心：18 世纪分析力学的发展完全建立在类似思想上。1851 年，Sylvester 试图构造线性群 \(\mathbf{GL}(3,\mathbf{C})\) 及其某些子群的不变量，把其矩阵中出现的参数视为形如 \(\alpha_{j}dt\) 的“无穷小”增量；他把函数 \(f((z_{j}))\) 是不变量这一事实写成方程 \(f((z_{j}+\alpha_{j}(dt))=f((z_{j}))\)；这给出了关于 f 的线性偏微分方程 Xf=0，其中

\[
\mathrm{X} f = \sum_ {j} \alpha_ {j} \frac {\partial f}{\partial z _ {j}},\tag{1}
\]

其中 X 因而是一个微分算子，即“沿方向参数 \(\alpha_{j}\) 的导数”（{[}5{]}，第 3 卷，第 326、327 页）；Sylvester 似乎认为这是一个极其重要的一般原理，但此后似乎从未再回到这个问题。不久之后，Cayley（{[}6{]}，第 II 卷，第 164--178 页）以类似方式研究某些表示下 \(\mathbf{SL}(2,\mathbf{C})\) 的不变量，并证明它们是两个一阶偏微分方程 \(Xf=0\)、\(Yf=0\) 的解，其中 X 和 Y 如上由“无穷小”变换得到

\[
\left( \begin{array}{c c} 0 & 0 \\ d t & 0 \end{array} \right) \quad \text { and } \quad \left( \begin{array}{c c} 0 & d t \\ 0 & 0 \end{array} \right).
\]

用现代术语说，这表述为 X 和 Y 生成李代数 \(\mathfrak{sl}(2,\mathbf{C})\)；此外，Cayley 明确计算了括号 XY -- YX，并证明它也来自一个“无穷小”变换。

在他 1868 年关于运动群的论文 {[}7{]} 中，Jordan 自始至终使用“无穷小变换”这一概念，但完全是从几何观点出发的。他无疑提出了由无穷小变换“生成”的单参数群这一思想：对 Jordan 而言，它是通过“适当地重复”无穷小变换而得到的变换集合（同上，第 243 页）。Klein 和 Lie 在他们的论文中使用了同一表述“重复的无穷小变换” {[}1 b{]}，但上下文表明，他们所指的是一个微分系统的积分。如果他们考虑的单参数群由变换 \(x' = f(x, y, t)\)、\(y' = g(x, y, t)\) 组成，则相应的“无穷小变换”由下式给出：

\[
d x = p (x, y) d t, \quad d y = q (x, y) d t
\]

其中 \(p(x,y) = \frac{\partial f}{\partial t} (x,y,t_0), q(x,y) = \frac{\partial g}{\partial t} (x,y,t_0)\)，且 \(t_0\) 对应于

该群的恒等变换。由于 Klein 和 Lie 明确知道函数 f 和 g，他们很容易验证下列函数

\[
t \mapsto f (x, y, t) \quad \text { and } \quad t \mapsto g (x, y, t)
\]

用参数形式给出微分方程的积分曲线

\[
q (\xi , \eta) d \xi = p (\xi , \eta) d \eta
\]

经过点 \((x,y)\)，但他们没有给出一般论证；而且在其论文余下部分也从未使用这一事实。

\textbf{(c) 接触变换}

接下来的两年中，Lie 似乎放弃了变换群理论（尽管他仍与于 1872 年发表“纲领”的 Klein 保持非常密切的联系），转而研究接触变换、一阶偏微分方程的积分以及这两个理论之间的关系。这里不讨论这些问题的历史，只提及几点似乎在变换群理论的起源中发挥了重要作用的事实。

接触变换的概念同时推广了点变换和逆极变换。粗略地说，\(\mathbf{C}^{\mathrm{n}}\) 上的接触变换† 是一个开子集 \(\Omega\) 的同构，其中该开子集属于流形 \(T^{\prime}(\mathbf{C}^{\mathrm{n}})\)，该流形由指向 \(\mathbf{C}^{\mathrm{n}}\) 的余切向量组成；另一个开子集 \(\Omega^{\prime}\) 属于 \(T^{\prime}(\mathbf{C}^{\mathrm{n}})\)，且该同构把 \(\Omega\) 的典范 1-形式映到 \(\Omega^{\prime}\) 的典范 1-形式。换言之，若 \((x_{1},\ldots ,x_{n},p_{1},\ldots ,p_{n})\) 是 \(T^{\prime}(\mathbf{C}^{\mathrm{n}})\) 的典范坐标，则接触变换是同构 \((x_{1},p_{1})\mapsto (\mathbf{X}_{\mathrm{i}},\mathbf{P}_{\mathrm{i}})\)，满足关系 \(\sum_{i = 1}^{n}\mathbf{P}_{i}d\mathbf{X}_{i} = \sum_{i = 1}^{n}p_{i}dx_{i}\)。这类变换出现在研究形如

\[
\mathrm{F} \left(x _ {1}, x _ {2}, \dots , x _ {n}, \frac {\partial z}{\partial x _ {1}}, \dots , \frac {\partial z}{\partial x _ {n}}\right) = 0.\tag{2}
\]

Lie 在研究这些问题的过程中熟悉了 Poisson 括号的运算

\[
(f, g) = \sum_ {i = 1} ^ {n} \left(\frac {\partial f}{\partial x _ {i}} \frac {\partial g}{\partial p _ {i}} - \frac {\partial g}{\partial x _ {i}} \frac {\partial f}{\partial p _ {i}}\right)\tag{3}
\]

以及微分算子 \(\ddagger_{\xi}\) 的括号 {[}X, Y{]} = XY - YX；他把 Poisson 括号 (3) 解释为与 g 对应的 (1) 型变换对 f 的作用，并注意到，Poisson 括号的 Jacobi 恒等式意味着，与 f 和 g 对应的微分算子的括号与括号 (g, h) 相联系。对满足 (F, g) = 0 的函数 g 的研究（它出现在 Jacobi 积分偏微分方程 (2) 的方法中）使 Lie 转而研究保持给定方程不变的无穷小接触变换。最后，Lie 研究了这样的函数集 \((u_{j})_{1 \in j \leq m}\)：它们是 \(x_{i}\) 和 \(p_{i}\) 的函数，并且括号 \((u_{j}, u_{k})\) 是 \(u_{h}\) 的函数；他称这些集合为“群”（它们实质上已由 Jacobi 研究过） {[}3 c{]}。

